\chapter{Bilan et perspectives}\label{chap:conclusion}

\section{Bilan}
Le projet livre une plateforme complète et opérationnelle : six services autonomes aux contrats explicites,
trois en SOAP décrits par WSDL et trois en REST, orchestrés par un processus BPMN 2.0 réellement exécuté.
Le processus couvre les neuf étapes demandées et va au-delà sur trois points : demande de complément d'information,
relance automatique de l'entreprise et interruption sur annulation. La gestion des erreurs combine reprise
automatique, événements d'erreur BPMN et traitement humain des incidents, avec des reprises idempotentes.
L'ensemble est vérifié par 28 tests d'intégration de bout en bout, sur H2 et PostgreSQL, et peut être utilisé par
chaque acteur à travers l'interface web.

\section{Limites}
\begin{itemize}
  \item Les six services partagent un processus et une base physique pour simplifier la démonstration ; l'isolation
        est logique (tables, contrats, appels réseau), pas physique.
  \item Le transport n'est pas chiffré dans l'environnement de démonstration (HTTPS à activer en production).
\end{itemize}

\section{Perspectives}
\begin{itemize}
  \item Déployer chaque service dans son propre conteneur avec sa base (Docker Compose ou Kubernetes) et un annuaire
        de services ou une passerelle d'API.
  \item Remplacer HTTP Basic par des jetons OAuth 2 / JWT, et signer les messages SOAP (WS-Security \emph{Signature}).
  \item Ajouter un disjoncteur (\emph{circuit breaker}) et une file de messages pour les notifications.
  \item Remplacer Mailpit par le serveur SMTP de l'école et ajouter des notifications instantanées dans l'application.
  \item Enrichir la compatibilité par une recommandation d'offres fondée sur l'historique des stages.
\end{itemize}
